% Candidate rules for airline-initiated change.  Written by hand; see docs/one-order-coverage.md.

\section{Airline-initiated change}

% Second survey round: candidates written from the survey notes in docs/one-order-survey.md.

\begin{theorem}[Penalty when declining an airline change (C-100)]
  \label{thm:airline-change-cancel-after-involuntary}
  \uses{thm:c16-penalty-net, def:c15-cancel-order, def:airline-change-follow-up-action, def:airline-change-involuntary-change}
  \notready
  A customer who declines an airline-initiated change through the Cancel
  follow-up action pays no cancellation penalty.

  Case to check: In the itinerary change worked example the new order item
  carries a 50 EUR cancellation fee. Applying it to a Cancel follow-up
  charges the customer for the airline's own change.

  Open point: The sources offer Cancel as a response to an airline change
  but never say whether penalties apply. The cancellation page covers only
  customer-initiated cancellation and leaves penalties to the airline.

  Sources: \href{https://standards.atlassian.net/wiki/spaces/EASD/pages/574849094}{EASD guide, Seller Follow-up Actions};
  \href{https://standards.atlassian.net/wiki/spaces/EASD/pages/672825400}{EASD guide, Inform of Impacted Dependent Services};
  \href{https://standards.atlassian.net/wiki/spaces/EASD/pages/574783569}{EASD guide, Cancellation of Order Items};
  \href{https://standards.atlassian.net/wiki/spaces/EASD/pages/815202352}{EASD guide, Airline notifies Seller of changes to an Order}.
\end{theorem}

\begin{theorem}[Refund for a cancelled flight (C-101)]
  \label{thm:airline-change-refund-undeliverable}
  \uses{def:c05-holds-value, def:c08-delivery-status, thm:c04-item-unsplit, def:refund, def:ancillaries-flight-change-dependents}
  \notready
  When the airline cancels a flight and does not re-accommodate, the flight
  service and its dependent services become \texttt{CANCELLED} with delivery
  status \texttt{UNABLE TO DELIVER}. The customer is then owed a refund of
  the value of those services.

  Case to check: A two-segment item costs 1000 EUR and the airline cancels
  the second segment. The item stays \texttt{ACTIVE} because the first
  segment still holds value, and no rule returns any part of the 1000 EUR.

  Open point: The worked example offers no follow-up action and no refund,
  and no source says whether \texttt{UNABLE TO DELIVER} caused by the
  airline entitles the customer to money back. Because price is set per
  order item, the business must also decide how much of an item's price
  belongs to one cancelled segment.

  Sources: \href{https://standards.atlassian.net/wiki/spaces/EASD/pages/815202352}{EASD guide, Airline notifies Seller of changes to an Order};
  \href{https://standards.atlassian.net/wiki/spaces/EASD/pages/574849062}{EASD guide, Order}.
\end{theorem}

\begin{definition}[Consequence of a missed time limit (C-102)]
  \label{def:airline-change-inaction-consequence}
  \uses{def:airline-change-follow-up-action, def:instant, def:c15-cancel-order, def:refund}
  \notready
  If a follow-up action carries a time limit and the seller has not acted by
  then, the airline applies the stated consequence. Auto-acceptance keeps
  the changed order and Auto-cancellation cancels it. The time limit and the
  consequence are both optional.

  Case to check: The concept page pairs Accept with Auto-cancellation. A
  seller who misses the deadline then costs the customer the whole trip
  after a change the airline made.

  Open point: The business must decide what applies when either the time
  limit or the consequence is missing (the Cancel example has neither), and
  what Auto-cancellation cancels and refunds. It must also decide whether
  the time limit has to fall before the affected flight departs.

  Sources: \href{https://standards.atlassian.net/wiki/spaces/EASD/pages/574849094}{EASD guide, Seller Follow-up Actions};
  \href{https://standards.atlassian.net/wiki/spaces/EASD/pages/815202352}{EASD guide, Airline notifies Seller of changes to an Order};
  \href{https://standards.atlassian.net/wiki/spaces/EASD/pages/672825400}{EASD guide, Inform of Impacted Dependent Services}.
\end{definition}

\begin{theorem}[Change applied before seller response (C-103)]
  \label{thm:airline-change-applied-before-acceptance}
  \uses{def:c06-item-status, thm:c21-version, thm:c19-transfer-balanced, def:airline-change-follow-up-action}
  \notready
  An airline-initiated change takes effect in the order when it is notified.
  The version rises, replaced services are cancelled or transferred, and the
  new order item is \texttt{ACTIVE} while the seller's Accept is still
  pending.

  Case to check: In the itinerary example the original item is already
  \texttt{CANCELLED} at version 2 while Accept is pending. A customer who
  rejects the change cannot get the original flights back, so Accept offers
  no real choice.

  Open point: This is shown only in a worked example. The business must
  decide whether a change awaiting acceptance is provisional, and whether
  declining it restores the original item or only lets the customer cancel
  or reshop from the changed order.

  Sources: \href{https://standards.atlassian.net/wiki/spaces/EASD/pages/815202352}{EASD guide, Airline notifies Seller of changes to an Order};
  \href{https://standards.atlassian.net/wiki/spaces/EASD/pages/574554168}{EASD guide, Seller Initiated Change to an Order}.
\end{theorem}

\begin{theorem}[Paid status after re-accommodation (C-104)]
  \label{thm:airline-change-reaccommodated-paid-status}
  \uses{def:c08-delivery-status, def:c11-payment-time-limit, thm:c19-transfer-balanced, thm:airline-change-applied-before-acceptance}
  \notready
  A service the airline creates to replace a paid service, with no price
  difference owed, starts with delivery status \texttt{READY TO PROCEED}.

  Case to check: If \texttt{CONFIRMED} is read as unpaid, a payment time
  limit could apply to a re-accommodated trip the customer already paid for
  in full.

  Open point: The itinerary worked example gives such services delivery
  status \texttt{CONFIRMED}, which the Order concept reserves for services
  still awaiting payment, although the exchange is even. The business must
  decide whether paid status carries across an airline-initiated change.

  Sources: \href{https://standards.atlassian.net/wiki/spaces/EASD/pages/815202352}{EASD guide, Airline notifies Seller of changes to an Order};
  \href{https://standards.atlassian.net/wiki/spaces/EASD/pages/574849062}{EASD guide, Order};
  \href{https://standards.atlassian.net/wiki/spaces/EASD/pages/672825400}{EASD guide, Inform of Impacted Dependent Services}.
\end{theorem}

\begin{theorem}[Unacknowledged change notice (C-105)]
  \label{thm:airline-change-lost-notice}
  \uses{def:airline-change-inaction-consequence, def:airline-change-notice-chain}
  \notready
  A follow-up consequence takes effect only after the seller has
  acknowledged the notice that carried it.

  Case to check: A notice carrying Accept with Auto-cancellation is lost.
  The seller sees the pending action only if it retrieves the order, and the
  order may be cancelled before it does.

  Open point: The guide names an unreceived or unacknowledged notice as the
  failure case but attaches no rule to it: no retry, no fallback, and no
  effect on the time limit. The business must decide whether the time limit
  runs from sending or from acknowledgement.

  Sources: \href{https://standards.atlassian.net/wiki/spaces/EASD/pages/574455903}{EASD guide, Notification of Airline Initiated Change to an Order};
  \href{https://standards.atlassian.net/wiki/spaces/EASD/pages/574849094}{EASD guide, Seller Follow-up Actions}.
\end{theorem}

\begin{definition}[Changes that require a customer choice (C-106)]
  \label{def:airline-change-change-needs-choice}
  \uses{def:airline-change-follow-up-action, def:airline-change-involuntary-change}
  \notready
  The airline may change departure times, flight numbers, aircraft,
  operating carrier or itinerary, or cancel a flight, and apply the change
  without the customer. For each change it decides whether to attach
  follow-up actions or send a notice only.

  Case to check: In the worked examples a 10-minute departure change and a
  change of operating carrier carry no follow-up action. A customer moved to
  another carrier has no recorded way to refuse.

  Open point: The sources give no threshold. The business must decide which
  change types and sizes oblige the airline to offer Accept, Cancel or
  Reshop.

  Sources: \href{https://standards.atlassian.net/wiki/spaces/EASD/pages/815202352}{EASD guide, Airline notifies Seller of changes to an Order};
  \href{https://standards.atlassian.net/wiki/spaces/EASD/pages/574455903}{EASD guide, Notification of Airline Initiated Change to an Order}.
\end{definition}

\begin{definition}[Voluntary and involuntary change (C-107)]
  \label{def:airline-change-involuntary-change}
  \uses{def:principal, def:order}
  \notready
  An order change is voluntary when the customer requests it and involuntary
  when the airline imposes it. Resolution 797 commits One Order to support
  servicing, delivery and accounting for both kinds.

  Open point: No source defines the two kinds. Change notices also carry
  changes the customer made directly with the airline, such as a name or
  contact update, so the business must decide which change reasons count as
  involuntary and trigger customer protections.

  Sources: \href{https://www.iata.org/contentassets/72cbd60393ff42b5975d90ce9e049a7d/oneorder-resolution-797.pdf}{IATA Resolution 797};
  \href{https://standards.atlassian.net/wiki/spaces/EASD/pages/672825400}{EASD guide, Inform of Impacted Dependent Services};
  \href{https://retailing.iata.org/armi/docs/ORDOC2/}{IATA capability ORDOC2};
  \href{https://standards.atlassian.net/wiki/spaces/EASD/pages/815202352}{EASD guide, Airline notifies Seller of changes to an Order}.
\end{definition}

\begin{theorem}[Pending follow-up actions (C-108)]
  \label{thm:airline-change-pending-actions}
  \uses{def:airline-change-follow-up-action, def:airline-change-inaction-consequence}
  \notready
  A follow-up action the seller has not performed appears in every order
  view until it is resolved. Performing any one of the actions offered for a
  change resolves all of them.

  Case to check: The itinerary example offers Accept, Cancel, Reshop and
  ContactAirline under one time limit with Auto-acceptance. The seller
  phones the airline, nothing in the order records it, and auto-acceptance
  fires.

  Open point: The source says pending actions are shown but not when they
  stop being pending. The business must decide whether one response clears
  the rest, and whether ContactAirline, which happens offline, ever clears a
  pending Accept.

  Sources: \href{https://standards.atlassian.net/wiki/spaces/EASD/pages/574849094}{EASD guide, Seller Follow-up Actions};
  \href{https://standards.atlassian.net/wiki/spaces/EASD/pages/815202352}{EASD guide, Airline notifies Seller of changes to an Order}.
\end{theorem}

\begin{definition}[Repricing after waitlist confirmation (C-109)]
  \label{def:airline-change-reprice-and-accept}
  \uses{def:airline-change-follow-up-action, def:airline-change-inaction-consequence, def:c08-delivery-status, def:c11-payment-time-limit}
  \notready
  When a waitlisted service is confirmed, the airline may ask the seller to
  reprice the order item and accept. The guide's example pairs this with
  Auto-acceptance if the seller does not act.

  Case to check: The waitlist clears at a higher fare. The seller misses the
  deadline, auto-acceptance confirms the service, and the customer is bound
  to a price never quoted.

  Open point: The source does not say at what price an auto-accepted
  waitlist confirmation stands, or whether the customer must pay a
  difference before the service becomes \texttt{READY TO PROCEED}.

  Source: \href{https://standards.atlassian.net/wiki/spaces/EASD/pages/574849094}{EASD guide, Seller Follow-up Actions}.
\end{definition}

\begin{definition}[Who is told of a change (C-110)]
  \label{def:airline-change-notice-chain}
  \uses{def:principal, def:order}
  \notready
  A service provider reports every service status change to the airline that
  owns the order. That airline notifies the seller of airline-initiated
  changes and the seller acknowledges. The seller is expected to inform the
  customer when a change needs action.

  Case to check: Agency A creates the order and agency B services it later.
  The disruption notice goes only to A, and B keeps selling changes against
  an old version.

  Open point: The sources do not say which seller is notified when several
  sellers have serviced the order, whether the airline must reach the
  customer when the seller does not acknowledge, or whether the seller's
  duty to inform the customer is binding.

  Sources: \href{https://standards.atlassian.net/wiki/spaces/EASD/pages/574455903}{EASD guide, Notification of Airline Initiated Change to an Order};
  \href{https://standards.atlassian.net/wiki/spaces/EASD/pages/672825400}{EASD guide, Inform of Impacted Dependent Services};
  \href{https://standards.atlassian.net/wiki/spaces/EASD/pages/574652557}{EASD guide, Order History};
  \href{https://retailing.iata.org/armi/docs/STPNTF/}{IATA capability STPNTF};
  \href{https://www.iata.org/contentassets/72cbd60393ff42b5975d90ce9e049a7d/oneorder-resolution-797.pdf}{IATA Resolution 797}.
\end{definition}

\begin{theorem}[Cancelled service in an active item (C-111)]
  \label{thm:airline-change-cancelled-in-active-item}
  \uses{def:c06-item-status, def:c08-delivery-status, def:c18-change}
  \notready
  A service with status \texttt{CANCELLED} belongs to an order item with
  status \texttt{CANCELLED}.

  Case to check: In the flight cancellation worked example the second
  segment's flight and bag services become \texttt{CANCELLED} while their
  order item gets no status change.

  Open point: The Order concept ties a \texttt{CANCELLED} service to a
  cancelled item, and the cancellation page treats cancelling some services
  of an item as a change. The flight cancellation example cancels one
  segment's services and leaves the item untouched, so the business must
  decide which model holds.

  Sources: \href{https://standards.atlassian.net/wiki/spaces/EASD/pages/574849062}{EASD guide, Order};
  \href{https://standards.atlassian.net/wiki/spaces/EASD/pages/574783569}{EASD guide, Cancellation of Order Items};
  \href{https://standards.atlassian.net/wiki/spaces/EASD/pages/815202352}{EASD guide, Airline notifies Seller of changes to an Order}.
\end{theorem}

\begin{theorem}[Price of reshopping after airline change (C-112)]
  \label{thm:airline-change-reshop-after-involuntary}
  \uses{def:c18-change, def:airline-change-follow-up-action, def:airline-change-involuntary-change}
  \notready
  A customer who declines an airline-initiated change through the Reshop
  follow-up action pays no change fee.

  Case to check: The new item in the itinerary example carries a 40 EUR
  change fee. A natural reading charges it when the customer reshops away
  from flights the airline moved.

  Open point: The guide says such a reshop may or may not lead to additional
  payment and gives no rule. The business must decide whether change fees
  and fare differences apply, and which alternatives the airline must offer.

  Sources: \href{https://standards.atlassian.net/wiki/spaces/EASD/pages/672825400}{EASD guide, Inform of Impacted Dependent Services};
  \href{https://standards.atlassian.net/wiki/spaces/EASD/pages/574849094}{EASD guide, Seller Follow-up Actions};
  \href{https://standards.atlassian.net/wiki/spaces/EASD/pages/815202352}{EASD guide, Airline notifies Seller of changes to an Order}.
\end{theorem}

\begin{definition}[Seller follow-up action (C-113)]
  \label{def:airline-change-follow-up-action}
  \uses{def:order, def:instant}
  A change notice carries zero or more seller follow-up actions drawn from
  Accept, Reshop, Cancel, ContactAirline and RepriceAndAccept. Each may
  carry a time limit, a consequence of inaction, and an airline contact
  link. A notice with no follow-up action needs only acknowledgement.

  Sources: \href{https://standards.atlassian.net/wiki/spaces/EASD/pages/574849094}{EASD guide, Seller Follow-up Actions};
  \href{https://standards.atlassian.net/wiki/spaces/EASD/pages/574455903}{EASD guide, Notification of Airline Initiated Change to an Order}.
\end{definition}
